\section{Introduction}
\begin{frame}{What is Unsupervised Learning?}
\textbf{So far:} We've been doing \textbf{Supervised Learning}, learning from labeled data $(x, y)$.

\vspace{1em}
\begin{columns}[c]
\begin{column}{0.50\textwidth}
\begin{tcolorbox}[colback=blue!5, colframe=blue!50, title=\textbf{Supervised Learning}]
\textbf{Data:} $(x, y)$ pairs
\begin{itemize}
\item $x$: input (features)
\item $y$: label (target)
\end{itemize}
\vspace{0.3em}
\textbf{Goal:} Learn $f(x) \approx y$

\textbf{Examples:}
\begin{itemize}
\item Image classification
\item House price prediction
\end{itemize}
\end{tcolorbox}
\end{column}

\begin{column}{0.50\textwidth}
\begin{tcolorbox}[colback=green!5, colframe=green!50!black, title=\textbf{Unsupervised Learning}]
\textbf{Data:} Only $x$ (no labels!)
\begin{itemize}
\item Just raw data
\item No target to predict
\end{itemize}
\vspace{0.3em}
\textbf{Goal:} Find hidden patterns or structure

\textbf{Examples:}
\begin{itemize}
\item Customer segmentation
\item Data compression
\end{itemize}
\end{tcolorbox}
\end{column}
\end{columns}

\vspace{1em}
\centering
\textbf{Key difference:} No labels $\rightarrow$ algorithm must discover structure on its own!
\end{frame}

% --- SLIDE: Why Unsupervised Learning? ---
\begin{frame}{Why Unsupervised Learning?}
\textbf{Labels are expensive!}

\vspace{0.5em}
\begin{itemize}
\item Labeling data requires human effort, time, and expertise
\item Most real-world data is unlabeled
\item Sometimes we don't even know what labels to look for
\end{itemize}

\vspace{1em}
\begin{tcolorbox}[colback=yellow!8, colframe=orange!60, title=\textbf{The Big Idea}]
Can we learn useful things from data \textbf{without} labels? \\
Yes! Unsupervised learning finds hidden patterns, structure, and representations.
\end{tcolorbox}

\end{frame}

% --- SLIDE: Tasks in Unsupervised Learning ---
\begin{frame}{Common Unsupervised Learning Tasks}
\textbf{What can we do without labels?}

% \vspace{0.8em}
\begin{columns}[t]
\begin{column}{0.52\textwidth}
\begin{tcolorbox}[colback=blue!5, colframe=blue!50, height=2.7cm]
\textbf{1. Clustering}

Group similar data points together

\textbf{Example:} Customer segmentation, document group
\end{tcolorbox}

\vspace{0.5em}

\begin{tcolorbox}[colback=red!5, colframe=red!50, height=2.7cm]
\textbf{2. Dimensionality Reduction}

Compress high-dimensional data to lower dimensions

\textbf{Example:} Visualization, feature extraction
\end{tcolorbox}
\end{column}

\begin{column}{0.50\textwidth}
\begin{tcolorbox}[colback=orange!5, colframe=orange!50, height=2.7cm]
\textbf{3. Anomaly Detection}

Identify unusual or outlier data points

\textbf{Example:} Fraud detection, quality control
\end{tcolorbox}

\vspace{0.5em}

\begin{tcolorbox}[colback=purple!5, colframe=purple!50, height=2.7cm]
\textbf{4. Generation}

Learn to create new, realistic data samples

\textbf{Example:} Image synthesis, data augmentation
\end{tcolorbox}
\end{column}
\end{columns}

\vspace{0.5em}
\centering
\small \textit{We'll explore four powerful algorithms that tackle these tasks.}
\end{frame}

% =========================
% K-Means Clustering - IMPROVED VERSION
% =========================

% --- SLIDE 1: What is Clustering? ---
